\documentclass[10pt,a4paper]{article}

\usepackage[utf8]{inputenc}
\usepackage[T1]{fontenc}
\usepackage{amsmath,amssymb,amsfonts}
\usepackage{mathtools}
\usepackage{bm}
\usepackage{booktabs}
\usepackage{hyperref}
\usepackage[margin=2.5cm]{geometry}
\usepackage{enumitem}
\usepackage{xcolor}

% Useful shortcuts
\newcommand{\Cl}{C_\ell}
\newcommand{\Nell}{N_\ell}
\newcommand{\Bell}{B_\ell}
\newcommand{\fsky}{f_{\mathrm{sky}}}
\newcommand{\Dl}{\Delta\ell}
\newcommand{\sigmav}{\langle\sigma v\rangle}
\newcommand{\mchi}{m_\chi}
\newcommand{\rhobar}{\bar{\rho}}
\newcommand{\dndm}{\frac{\mathrm{d}n}{\mathrm{d}M}}
\newcommand{\utilde}{\tilde{u}}
\newcommand{\vtilde}{\tilde{v}}
\newcommand{\nhat}{\hat{\mathbf{n}}}

\title{Statistical Framework for HI $\times$ Gamma-Ray \\
Cross-Correlation Power Spectrum Analysis \\
in Dark Matter Searches}
\author{}
\date{}

\begin{document}
% \maketitle

% \begin{abstract}
% The Gaussian covariance of the cross-correlation angular power spectrum between HI 21-cm intensity maps and Fermi-LAT gamma-ray maps is the statistical backbone enabling dark matter searches through this novel multi-messenger channel. Pinetti, Camera, Fornengo \& Regis (2020) demonstrated that MeerKAT $\times$ Fermi-LAT can achieve a $3.7\sigma$ detection of the astrophysical cross-correlation signal, SKA Phase~1 reaches $5.7\sigma$, and SKA Phase~2 achieves $8.2\sigma$, while probing thermal-relic WIMP annihilation up to $\mchi \approx 130$~GeV. The statistical framework rests on a chain of results: Wick's theorem decomposing Gaussian four-point functions, the Knox (1995) formula incorporating instrumental noise and beam, pseudo-$\Cl$ estimators correcting for incomplete sky coverage, and hypothesis testing via $\Delta\chi^2$ raster scans or profile likelihood methods. This report derives these results from first principles, analyzes their assumptions and limitations, and extends the discussion to non-Gaussian corrections, practical estimators, and advanced inference frameworks.
% \end{abstract}

% \tableofcontents
% \newpage

%======================================================================
\section{Deriving the Gaussian Covariance from Wick's Theorem}
\label{sec:covariance}
%======================================================================

\subsection{Spherical Harmonic Setup and the Power Spectrum Estimator}

Any sky field (gamma-ray intensity or HI brightness temperature) is expanded in spherical harmonics as
\begin{equation}
T(\nhat) = \sum_{\ell m} a_{\ell m}\, Y_\ell^m(\nhat).
\end{equation}
For statistically isotropic Gaussian random fields, the two-point function of harmonic coefficients between fields $i$ and $j$ satisfies
\begin{equation}
\langle a_{\ell m}^i \, a_{\ell' m'}^{j*} \rangle = \Cl^{ij}\, \delta_{\ell\ell'}\, \delta_{mm'},
\end{equation}
where $\Cl^{ij}$ is the angular cross-power spectrum. The natural unbiased estimator is
\begin{equation}
\hat{C}_\ell^{ij} = \frac{1}{2\ell+1} \sum_{m=-\ell}^{+\ell} a_{\ell m}^i \, a_{\ell m}^{j*}.
\label{eq:cl_estimator}
\end{equation}
Computing the covariance of this estimator requires evaluating the four-point function
$\langle a_{\ell m}^i \, a_{\ell m}^{j*} \, a_{\ell' m'}^k \, a_{\ell' m'}^{n*} \rangle$.


\subsection{Wick's Theorem Decomposition into Three Pairings}

\textbf{Wick's theorem} (equivalently, Isserlis' theorem for real Gaussian variables) states that the expectation of a product of $2n$ zero-mean jointly Gaussian variables decomposes into a sum over all distinct pairings. For four variables, exactly \textbf{three pairings} arise:
\begin{equation}
\langle X_1 X_2 X_3 X_4 \rangle
= \langle X_1 X_2 \rangle \langle X_3 X_4 \rangle
+ \langle X_1 X_3 \rangle \langle X_2 X_4 \rangle
+ \langle X_1 X_4 \rangle \langle X_2 X_3 \rangle.
\end{equation}

Identifying $X_1 = a_{\ell m}^i$, $X_2 = a_{\ell m}^{j*}$, $X_3 = a_{\ell' m'}^k$, $X_4 = a_{\ell' m'}^{n*}$:

\begin{itemize}[leftmargin=2em]
\item \textbf{Pairing~1}: $\langle a_{\ell m}^i\, a_{\ell m}^{j*}\rangle\, \langle a_{\ell' m'}^k\, a_{\ell' m'}^{n*}\rangle = \Cl^{ij} \cdot C_{\ell'}^{kn}$. This is the disconnected piece that cancels when computing the covariance (it equals $\langle \hat{C}_\ell^{ij}\rangle \langle \hat{C}_{\ell'}^{kn}\rangle$).

\item \textbf{Pairing~2}: $\langle a_{\ell m}^i\, a_{\ell' m'}^{n*}\rangle\, \langle a_{\ell m}^{j*}\, a_{\ell' m'}^k\rangle = \Cl^{in}\, \Cl^{jk}\, \delta_{\ell\ell'}\, \delta_{mm'}$. The Kronecker deltas force $\ell = \ell'$ and $m = m'$.

\item \textbf{Pairing~3}: $\langle a_{\ell m}^i\, a_{\ell' m'}^k\rangle\, \langle a_{\ell m}^{j*}\, a_{\ell' m'}^{n*}\rangle$. Using the reality condition ($a_{\ell,-m} = (-1)^m a_{\ell m}^*$) on both unconjugated factors, each two-point function picks up a phase $(-1)^{m'}$, giving a net factor $(-1)^{2m'} = 1$. The result is $\Cl^{ik}\, \Cl^{jn}\, \delta_{\ell\ell'}\, \delta_{m,-m'}$.
\end{itemize}

When summing over $m$ and $m'$, Pairing~2 contributes $(2\ell+1)$ terms (from $\delta_{mm'}$), and Pairing~3 also contributes $(2\ell+1)$ terms (from $\delta_{m,-m'}$, with the net phase $(-1)^{2m'} = 1$). After dividing by $(2\ell+1)^2$, the \textbf{full-sky Gaussian covariance} is:
\begin{equation}
\boxed{
\mathrm{Cov}\!\left(\hat{C}_\ell^{ij},\, \hat{C}_{\ell'}^{kn}\right)
= \frac{\delta_{\ell\ell'}}{2\ell+1}
\left[ \Cl^{ik}\, \Cl^{jn} + \Cl^{in}\, \Cl^{jk} \right].
}
\label{eq:gaussian_cov_fullsky}
\end{equation}

The factor $\mathbf{1/(2\ell+1)}$ reflects the number of independent $m$-modes available at each multipole---the fundamental origin of cosmic variance. At low~$\ell$, few modes exist and variance is large; at high~$\ell$, many modes are averaged and variance shrinks.


\subsection{Specializing to the Cross-Spectrum Variance}

For the cross-spectrum between fields $i$ (gamma-ray) and $j$ (HI), setting $k = i$ and $n = j$:
\begin{equation}
\mathrm{Var}\!\left(\hat{C}_\ell^{ij}\right)
= \frac{1}{2\ell+1}
\left[ \left(\Cl^{ij}\right)^2 + \Cl^{ii}\, \Cl^{jj} \right].
\label{eq:cross_var}
\end{equation}

The first term $(\Cl^{ij})^2$ captures the ``signal $\times$ signal'' cosmic variance of the cross-correlation itself, while the second term $\Cl^{ii}\, \Cl^{jj}$---the product of auto-power spectra---typically \textbf{dominates} because the cross-correlation coefficient
\begin{equation}
r_\ell = \frac{\Cl^{ij}}{\sqrt{\Cl^{ii}\, \Cl^{jj}}}
\end{equation}
is much less than unity for HI~$\times$~gamma-ray.


\subsection{Including Noise, Beam, and Partial Sky: Pinetti et al.\ Eq.~2.7}

In practice, each field is observed through an instrumental beam $\Bell^i$ and contaminated by noise with power spectrum $\Nell^i$. The observed harmonic coefficients are
\begin{equation}
d_{\ell m}^i = \Bell^i \, a_{\ell m}^i + n_{\ell m}^i,
\end{equation}
where noise between different experiments is \textbf{uncorrelated} ($\langle n^i\, n^{j*}\rangle = 0$ for $i \neq j$). After deconvolving beams, the effective auto-spectrum entering the covariance becomes
\begin{equation}
\tilde{C}_\ell^{ii} = \Cl^{ii} + \frac{\Nell^i}{(\Bell^i)^2}.
\end{equation}
For partial sky coverage with fraction $\fsky$, independent modes are reduced by a factor $\fsky$. With bandpower binning of width $\Dl$, Pinetti et al.\ Eq.~2.7 reads:
\begin{equation}
\boxed{
\mathrm{Var}\!\left(\hat{C}_\ell^{ij}\right)
= \frac{1}{(2\ell+1)\, \fsky\, \Dl}
\left[ \left(\Cl^{ij}\right)^2
+ \left(\Cl^{ii} + \frac{\Nell^i}{(\Bell^i)^2}\right)
  \left(\Cl^{jj} + \frac{\Nell^j}{(\Bell^j)^2}\right)
\right].
}
\label{eq:pinetti_2.7}
\end{equation}

This extends naturally to multiple energy bins and redshift bins, with the covariance becoming a matrix over these indices.


\subsection{Anatomy of Each Ingredient}

\paragraph{The cross-correlation term $(\Cl^{ij})^2$.}
This term is subdominant. The cross-correlation coefficient between the unresolved gamma-ray background and HI intensity is small ($r_\ell \ll 1$), because the gamma-ray sky is dominated by unresolved blazars and other sources that do not perfectly trace the HI distribution. This term can typically be neglected in forecasts. It would become non-negligible only in the (favorable) scenario where the cross-correlation signal is very strong---which would itself constitute a robust detection.

\paragraph{The gamma-ray auto-spectrum plus noise.}
This is $\tilde{C}_\ell^{\gamma\gamma} = \Cl^{\gamma\gamma} + N_\ell^\gamma / (B_\ell^\gamma)^2$. For a photon-count map with $N_\gamma$ detected photons over sky fraction $\fsky$, the Poisson noise angular power spectrum of the fractional overdensity field is $C_N = 4\pi \fsky / N_\gamma$ (white, scale-independent). When working with intensity maps (as in Pinetti et al.\ Table~2), the noise acquires intensity units: $N_\ell^\gamma = \bar{I}^2 \times 4\pi \fsky / N_\gamma$, where $\bar{I}$ is the mean intensity. In either convention, photon noise \textbf{dominates at all multipoles} for Fermi-LAT. The Fermi-LAT beam (point spread function) is energy-dependent:
\begin{equation}
B_\ell^\gamma = \exp\!\left[-\frac{\ell(\ell+1)\,\sigma_{\mathrm{beam}}^2}{2}\right],
\end{equation}
with the 68\% containment angle ranging from $\sim 5^\circ$ at 100~MeV to $\sim 0.8^\circ$ at 1~GeV and $\sim 0.2^\circ$ at 10~GeV (P8R3 SOURCE class). At high~$\ell$, $B_\ell^\gamma \to 0$, causing $N_\ell^\gamma/(B_\ell^\gamma)^2 \to \infty$ and killing the signal-to-noise.

\paragraph{The HI auto-spectrum plus noise.}
This is $\tilde{C}_\ell^{\mathrm{HI,HI}} = \Cl^{\mathrm{HI,HI}} + N_\ell^{\mathrm{HI}}/(B_\ell^{\mathrm{HI}})^2$. For \textbf{MeerKAT}, thermal noise dominates:
\begin{equation}
N_\ell^{\mathrm{HI}} \propto \frac{T_{\mathrm{sys}}^2}
{t_{\mathrm{pix}}\, \Delta\nu\, N_{\mathrm{dish}}},
\end{equation}
where $T_{\mathrm{sys}} \approx 25$--$30$~K. This places MeerKAT firmly in the \textbf{noise-dominated regime} where $N_\ell^{\mathrm{HI}} \gg \Cl^{\mathrm{HI,HI}}$. For \textbf{SKA Phase~2}, the dramatically improved sensitivity pushes the HI auto-spectrum toward the \textbf{cosmic-variance-dominated regime} where $\Cl^{\mathrm{HI,HI}} \gg N_\ell^{\mathrm{HI}}$ at low and intermediate multipoles. In the noise-dominated regime, more observation time helps linearly ($N_\ell \propto 1/t_{\mathrm{obs}}$); in the cosmic-variance-dominated regime, only surveying more sky area helps (variance $\propto 1/\fsky$).

\paragraph{Handling different sky coverage.}
When the gamma-ray survey (Fermi-LAT, nearly all-sky) and the radio survey (MeerKAT, $\sim 4{,}000$~deg$^2$) have different masks $W_\gamma(\nhat)$ and $W_{\mathrm{HI}}(\nhat)$, the effective sky fraction for the cross-correlation is determined by the \textbf{overlap region}. In the simplest approximation:
\begin{equation}
\fsky^{ij} = \frac{1}{4\pi}\int W_\gamma(\nhat)\, W_{\mathrm{HI}}(\nhat)\, \mathrm{d}\Omega.
\end{equation}
The more rigorous prescription (Fosalba et al.\ 2007) uses
\begin{equation}
\fsky = \frac{1}{4\pi}\int W_\gamma^2(\nhat)\, W_{\mathrm{HI}}^2(\nhat)\, \mathrm{d}\Omega
\end{equation}
for the covariance, accounting for the quadratic dependence on masks in the four-point function.

\paragraph{The beam-created ``sweet spot.''}
At low~$\ell$, cosmic variance dominates (few independent modes). At high~$\ell$, the Fermi-LAT PSF exponentially suppresses the gamma-ray signal while $N_\ell^\gamma/(B_\ell^\gamma)^2$ explodes. The cross-correlation SNR peaks at an \textbf{intermediate multipole range}---typically $\ell \sim 100$--$500$---where noise is manageable and enough modes exist for statistical power.


%======================================================================
\section{Combining Multipoles, Energy Bins, and Redshift Bins}
\label{sec:snr}
%======================================================================

\subsection{The SNR Formula (Pinetti et al.\ Eq.~5.6)}

The cumulative signal-to-noise ratio for detecting the cross-correlation signal aggregates information across all available dimensions. In the most general form:
\begin{equation}
\left(\frac{S}{N}\right)^2
= \sum_\ell \sum_{i,j} \sum_{a,b}
C_\ell^{\gamma_i,\, \mathrm{HI}_a}\;
\left[\mathrm{Cov}^{-1}\right]_{(i\ell a),(j\ell b)}\;
C_\ell^{\gamma_j,\, \mathrm{HI}_b},
\label{eq:snr_general}
\end{equation}
where the sums run over multipoles $\ell$ (from $\ell_{\min}$ to $\ell_{\max}$), gamma-ray energy bins $i,j$, and HI redshift bins $a,b$. In the simplified case of a single energy and redshift bin with diagonal covariance:
\begin{equation}
\left(\frac{S}{N}\right)^2
= \sum_\ell (2\ell+1)\, \fsky\, \Dl
\times \frac{(\Cl^{\gamma,\mathrm{HI}})^2}
{(\Cl^{\gamma,\mathrm{HI}})^2 + \tilde{C}_\ell^{\gamma\gamma}\, \tilde{C}_\ell^{\mathrm{HI,HI}}}.
\label{eq:snr_simple}
\end{equation}


\subsection{Independence Assumptions and Their Validity}

\paragraph{Independence across multipoles.}
The Gaussian covariance is diagonal in~$\ell$ ($\delta_{\ell\ell'}$), making different multipoles statistically independent. On a cut sky, the mode-coupling matrix $M_{\ell\ell'}$ introduces off-diagonal correlations, but after pseudo-$\Cl$ decoupling and bandpower binning, residual correlations are typically small. The $\fsky$ correction approximately accounts for the lost modes. For surveys covering $\fsky \gtrsim 0.05$, this approximation is accurate to $\sim 10$--$25\%$ (Challinor \& Chon 2005).

\paragraph{Independence across energy bins.}
This assumes that gamma-ray source populations contribute independently to different energy bands, or equivalently that the covariance matrix is block-diagonal in energy. In reality, the same astrophysical sources (blazars, SFGs) contribute correlated signals across energy bins. The Pinetti et al.\ framework handles this through the full covariance matrix $\mathrm{Cov}(\Cl^{\gamma_i,\mathrm{HI}_a},\, \Cl^{\gamma_j,\mathrm{HI}_b})$, which couples energy bins $i$ and $j$ through the gamma-ray auto-spectrum $\Cl^{\gamma_i,\gamma_j}$. Treating energy bins as independent (using only the diagonal $i=j$ elements) \textbf{overestimates} the SNR: positive off-diagonal covariance elements produce negative off-diagonal entries in the inverse covariance, which reduce the quadratic form $\mathbf{C}^\top \bm{\Sigma}^{-1}\mathbf{C}$ relative to the diagonal approximation. Neglecting these correlations is therefore \textbf{optimistic}, not conservative. In practice, however, photon noise is uncorrelated across energy bins ($N_\ell^{\gamma_i,\gamma_j} = 0$ for $i \neq j$), and because photon noise dominates the gamma-ray auto-spectrum, the off-diagonal covariance is suppressed and the diagonal approximation remains reasonable.

\paragraph{Independence across redshift bins.}
The HI intensity mapping survey provides precise redshift information through frequency channelization (the 21-cm line frequency maps directly to redshift). Redshift bins are independent when their separation exceeds the radial correlation length, which is typically $\Delta z \gtrsim 0.1$ for the relevant structures. For finer tomography, the correlation between neighboring bins is controlled by the overlap of their Limber kernels. The cross-power spectrum between redshift bins $a$ and $b$ is:
\begin{equation}
C_\ell^{ab} \propto
\int \frac{\mathrm{d}\chi}{\chi^2}\;
W_a(\chi)\, W_b(\chi)\, P(k{=}\ell/\chi),
\end{equation}
which vanishes when the window functions $W_a$ and $W_b$ do not overlap. Since this cross-spectrum enters the off-diagonal elements of the covariance matrix (through the Gaussian formula Eq.~\ref{eq:gaussian_cov_fullsky}), well-separated bins with $C_\ell^{ab} \approx 0$ are effectively independent.


\subsection{Multipole Range and Bandpower Binning}

The analysis adopts $\ell_{\min} = 10$ (set by the survey footprint: $\ell_{\min} \sim \pi/\theta_{\max}$) and $\ell_{\max} = 1000$--$2000$ (set by the worse angular resolution between the two instruments---typically the Fermi-LAT PSF). The exact $\ell_{\max}$ depends on energy: at $E \sim 1$~GeV, the Fermi-LAT PSF limits useful information to $\ell \lesssim 200$, while at $E \sim 10$~GeV, $\ell_{\max} \sim 1000$--$2000$ becomes accessible.

Bandpower binning groups multipoles into bins of width $\Dl$, serving two purposes: regularizing the mode-coupling matrix inversion and reducing the dimensionality of the data vector. The SNR is preserved under binning as long as the signal does not vary rapidly within a bin. Typical choices are linear bins with $\Dl = 20$--$50$ or logarithmic bins. The $(2\ell+1)\Dl$ factor in the SNR formula accounts for the number of modes within each bandpower.

For bandpowers, the SNR formula~(\ref{eq:snr_simple}) becomes:
\begin{equation}
\left(\frac{S}{N}\right)^2
= \sum_b \nu_b
\times \frac{(\bar{C}_b^{\gamma,\mathrm{HI}})^2}
{(\bar{C}_b^{\gamma,\mathrm{HI}})^2 + \bar{\tilde{C}}_b^{\gamma\gamma}\, \bar{\tilde{C}}_b^{\mathrm{HI,HI}}},
\label{eq:snr_bandpower}
\end{equation}
where $\nu_b = \sum_{\ell \in b}(2\ell+1)\fsky$ is the effective number of modes in bandpower~$b$ and $\bar{C}_b$ denotes the bandpower-averaged spectrum.


%======================================================================
\section{The $\Delta\chi^2$ Test for Dark Matter Detectability}
\label{sec:deltachi2}
%======================================================================

\subsection{Null Versus Alternative Hypotheses}

The $\Delta\chi^2$ framework (Pinetti et al.\ Eq.~5.7) tests whether a dark matter annihilation signal is detectable above the astrophysical background:
\begin{itemize}[leftmargin=2em]
\item \textbf{Null hypothesis $H_0$}: $\Cl = \Cl^{\mathrm{astro}}$ (blazars + star-forming galaxies + misaligned AGN only).
\item \textbf{Alternative hypothesis $H_1$}: $\Cl = \Cl^{\mathrm{astro}} + \Cl^{\mathrm{DM}}(\mchi, \sigmav)$.
\end{itemize}

The test statistic is:
\begin{equation}
\Delta\chi^2
= \sum_\ell \sum_{i,j} \sum_{a,b}
C_\ell^{\mathrm{DM},\gamma_i,\mathrm{HI}_a}\;
\left[\mathrm{Cov}^{-1}\right]_{(i\ell a),(j\ell b)}\;
C_\ell^{\mathrm{DM},\gamma_j,\mathrm{HI}_b},
\label{eq:deltachi2}
\end{equation}
where $\Cl^{\mathrm{DM}}$ is the DM-only contribution to the cross-power spectrum. Since $\Cl^{\mathrm{DM}} \propto \sigmav$ at fixed $\mchi$ and annihilation channel, the $\Delta\chi^2$ scales as $\sigmav^2$. The \textbf{$2\sigma$ (95.45\% CL) upper limit} on $\sigmav$ at fixed $\mchi$ is obtained by requiring:
\begin{equation}
\Delta\chi^2(\sigmav_{95}) = 4 \qquad \text{($2\sigma$ for 1 degree of freedom)}.
\end{equation}
Note that $\Delta\chi^2 = 3.84$ gives the exact 95\% CL; $\Delta\chi^2 = 4$ (i.e., $n_\sigma^2 = 2^2$) is standard practice in the particle physics literature (Pinetti et al.\ adopt this convention). The $5\sigma$ detection threshold corresponds to $\Delta\chi^2 = 25$.


\subsection{Raster Scan Methodology}

The raster scan fixes $\mchi$ on a grid spanning $\sim 1$~GeV to $\sim 10$~TeV, and at each mass point scans $\sigmav$ to find the excluded value. Since only $\sigmav$ varies while $\mchi$ is fixed, this is a \textbf{1-dof test}, and Wilks' theorem guarantees the asymptotic $\chi^2(1)$ distribution of $-2\ln(L_0/L_1)$. The 95\% exclusion curve $\sigmav_{95}(\mchi)$ traces the boundary in the $(\mchi, \sigmav)$ plane.

\textbf{Key results from Pinetti et al.\ Figure~12}: Fermi-LAT $\times$ SKA can probe the thermal relic cross-section $\sigmav \approx 3 \times 10^{-26}$~cm$^3$~s$^{-1}$ up to $\mchi \approx 130$~GeV for the $b\bar{b}$ channel. A future enhanced gamma-ray detector combined with an improved radio telescope extends this to the full WIMP mass window up to the TeV scale.


\subsection{Variance Under Null Versus Alternative}

A critical subtlety: the covariance matrix in the $\Delta\chi^2$ formula should strictly be evaluated under the \textbf{null hypothesis} (no DM signal) for computing \textbf{expected sensitivity} (the Asimov dataset approach of Cowan, Cranmer, Gross \& Vitells 2011). Using the null covariance gives the \textbf{median expected exclusion}---the limit one would obtain in 50\% of experimental realizations. Evaluating the covariance under $H_1$ (including the DM signal) is more conservative, as the additional DM contribution inflates the variance. The difference matters primarily when the DM signal significantly alters the total power spectrum, which occurs near the thermal relic cross-section for favorable mass ranges.

Explicitly, the covariance under $H_1$ replaces $\Cl^{\gamma\gamma}$ with $\Cl^{\gamma\gamma} + \Cl^{\gamma,\mathrm{DM}} + \Cl^{\mathrm{DM,DM}}$ in the auto-spectrum term. For weak signals ($\Cl^{\mathrm{DM}} \ll \Cl^{\gamma\gamma}$), the two approaches give identical results.


\subsection{Astrophysical Model Uncertainties}

The astrophysical cross-correlation model depends on luminosity functions of unresolved gamma-ray sources, their spectral energy distributions, and the HI halo occupation distribution. Uncertainty in these models propagates directly into the sensitivity forecast, because the covariance matrix contains the astrophysical auto-spectra. Marginalization over astrophysical parameters can be incorporated by either:
\begin{enumerate}[leftmargin=2em]
\item Profiling over nuisance parameters in a likelihood framework (Section~\ref{sec:advanced}).
\item Analytically marginalizing in a Fisher matrix approach.
\item Including Gaussian priors on the astrophysical amplitudes in the $\chi^2$.
\end{enumerate}

The tomographic information from redshift binning is crucial for breaking degeneracies between DM and astrophysical components, since their redshift evolution differs fundamentally: $\propto \rho_{\mathrm{DM}}^2(z)$ for DM versus the redshift-dependent luminosity density for astrophysical sources.


%======================================================================
\section{Beyond $\Delta\chi^2$: Sophisticated Inference Frameworks}
\label{sec:advanced}
%======================================================================

\subsection{Profile Likelihood with Nuisance Parameters}

The profile likelihood ratio provides a more rigorous treatment than the simple $\Delta\chi^2$:
\begin{equation}
\lambda(\mu) = \frac{L(\mu,\, \hat{\hat{\bm{\theta}}}(\mu))}
{L(\hat{\mu},\, \hat{\bm{\theta}})},
\label{eq:profile_lr}
\end{equation}
where $\mu = \sigmav / \sigmav_{\mathrm{ref}}$ is the signal strength, $\hat{\hat{\bm{\theta}}}(\mu)$ are the conditional MLEs of nuisance parameters $\bm{\theta}$ (astrophysical model parameters, systematic uncertainties) at fixed $\mu$, and $(\hat{\mu}, \hat{\bm{\theta}})$ are the unconditional MLEs. The test statistic $q_\mu = -2\ln\lambda(\mu)$ follows a $\chi^2(1)$ distribution asymptotically (Wald approximation; Cowan et al.\ 2011). When $\mu = 0$ lies at a physical boundary ($\sigmav \geq 0$), the distribution becomes a mixture:
\begin{equation}
f(q_0) = \tfrac{1}{2}\,\delta(q_0) + \tfrac{1}{2}\,\chi^2_1(q_0).
\end{equation}

The \textbf{CL$_s$ method} provides conservative upper limits by computing
\begin{equation}
\mathrm{CL}_s(\mu) = \frac{p_{s+b}(\mu)}{\mathrm{CL}_b(\mu)},
\end{equation}
preventing exclusion of signals the experiment has no sensitivity to---crucial when downward background fluctuations could produce artifactually strong limits. This is the standard approach for Fermi-LAT dwarf spheroidal galaxy analyses.


\subsection{Fisher Matrix Forecasts}

The Fisher information matrix for angular cross-power spectra between two fields is:
\begin{equation}
F_{\alpha\beta}
= \sum_\ell \frac{(2\ell+1)\fsky}{2}\;
\mathrm{Tr}\!\left[
\mathbf{C}_\ell^{-1}\, \frac{\partial \mathbf{C}_\ell}{\partial \theta_\alpha}\;
\mathbf{C}_\ell^{-1}\, \frac{\partial \mathbf{C}_\ell}{\partial \theta_\beta}
\right],
\label{eq:fisher}
\end{equation}
where $\mathbf{C}_\ell$ is the $2\times 2$ covariance matrix containing auto- and cross-spectra (plus noise):
\begin{equation}
\mathbf{C}_\ell = \begin{pmatrix}
\Cl^{\gamma\gamma} + N_\ell^\gamma/(B_\ell^\gamma)^2 & \Cl^{\gamma,\mathrm{HI}} \\
\Cl^{\gamma,\mathrm{HI}} & \Cl^{\mathrm{HI,HI}} + N_\ell^{\mathrm{HI}}/(B_\ell^{\mathrm{HI}})^2
\end{pmatrix}.
\end{equation}
The Cram\'er--Rao bound $\sigma(\theta_\alpha) \geq \sqrt{(F^{-1})_{\alpha\alpha}}$ sets the theoretical floor on parameter estimation. Marginalization over nuisance parameters via the Schur complement
\begin{equation}
F_{\mathrm{marg}} = F_{pp} - F_{pn}\, F_{nn}^{-1}\, F_{np}
\end{equation}
quantifies how much sensitivity is lost to astrophysical model uncertainty. Camera et al.\ (2015) pioneered the tomographic-spectral Fisher approach for gamma-ray $\times$ LSS cross-correlations, demonstrating that combined energy + redshift binning dramatically improves DM sensitivity.


\subsection{Template Fitting as Linear Regression}

The observed cross-power spectrum is decomposed as:
\begin{equation}
\Cl^{\mathrm{obs}}(E, z)
= \sum_\alpha A_\alpha\, T_\ell^\alpha(E, z),
\label{eq:template_fit}
\end{equation}
where $T_\ell^\alpha$ are templates for each source class $\alpha$ (BL~Lacs, FSRQs, SFGs, mAGN, DM for each annihilation channel). Each template has a \textbf{known shape} in $(\ell, E, z)$ space and an unknown amplitude $A_\alpha$. The amplitude vector is obtained via weighted least squares:
\begin{equation}
\hat{\mathbf{A}} = \left(\mathbf{T}^\top \bm{\Sigma}^{-1} \mathbf{T}\right)^{-1}
\mathbf{T}^\top \bm{\Sigma}^{-1} \mathbf{d},
\label{eq:wls}
\end{equation}
with covariance on the amplitudes:
\begin{equation}
\mathrm{Cov}(\hat{\mathbf{A}}) = \left(\mathbf{T}^\top \bm{\Sigma}^{-1} \mathbf{T}\right)^{-1}.
\label{eq:wls_cov}
\end{equation}
Here $\bm{\Sigma}$ is the full data covariance matrix and $\mathbf{d}$ is the data vector of measured bandpowers.

This exploits two orthogonal discriminants:
\begin{itemize}[leftmargin=2em]
\item \textbf{Spectral shape}: DM templates have a characteristic cutoff at $E = \mchi$, while blazars follow power laws.
\item \textbf{Angular shape}: DM one-halo term traces the NFW profile Fourier transform $\vtilde(k|M)$, while unresolved point sources produce Poissonian (flat) $\Cl$.
\end{itemize}

The condition number of $\mathbf{T}^\top \bm{\Sigma}^{-1} \mathbf{T}$ determines how well the decomposition can distinguish source classes. If templates are nearly degenerate (e.g., mAGN and blazars with similar spectra), the matrix becomes ill-conditioned and individual amplitudes are poorly constrained, even though the total signal is detected.


\subsection{Bayesian Inference and Model Comparison}

The Bayesian posterior is
\begin{equation}
P(\bm{\theta}|\mathbf{d}) \propto L(\mathbf{d}|\bm{\theta})\, \pi(\bm{\theta}),
\end{equation}
where $\pi(\bm{\theta})$ is the prior. Standard priors for DM searches include log-uniform distributions on $\sigmav$ (spanning $10^{-28}$--$10^{-22}$~cm$^3$~s$^{-1}$) and $\mchi$ (1--$10^4$~GeV), reflecting ignorance across orders of magnitude.

\textbf{Bayesian model comparison} via the Bayes factor
\begin{equation}
B_{10} = \frac{Z_1}{Z_0}
= \frac{\int L(\mathbf{d}|\bm{\theta}_1)\, \pi(\bm{\theta}_1)\, \mathrm{d}\bm{\theta}_1}
{\int L(\mathbf{d}|\bm{\theta}_0)\, \pi(\bm{\theta}_0)\, \mathrm{d}\bm{\theta}_0}
\end{equation}
naturally implements Occam's razor: the DM model is penalized for its extra parameters unless the data strongly favor them. Evidence computation via nested sampling (MultiNest, PolyChord) handles the multi-dimensional integral efficiently. On the modified Jeffreys scale (Trotta 2008), $|\ln B| > 5$ constitutes strong evidence.


%======================================================================
\section{Non-Gaussian Corrections to the Covariance}
\label{sec:nongaussian}
%======================================================================

\subsection{Physical Origins}

The full covariance decomposes as:
\begin{equation}
\mathrm{Cov} = \mathrm{Cov}^{\mathrm{G}} + \mathrm{Cov}^{\mathrm{cNG}} + \mathrm{Cov}^{\mathrm{SSC}}.
\end{equation}
The non-Gaussian terms arise from the connected trispectrum $T_{\ell_1\ell_2\ell_3\ell_4}$, generated by nonlinear gravitational evolution even from Gaussian initial conditions. In the halo model, the trispectrum separates into 1-halo, 2-halo, 3-halo, and 4-halo contributions, with the \textbf{1-halo term} dominating on small scales.


\subsection{Super-Sample Covariance}

Super-sample covariance (SSC) arises from coupling of sub-survey modes to super-survey (background) density fluctuations. Long-wavelength modes modulate the effective background density within the finite survey volume, coherently shifting the amplitude of all measured power spectra. The SSC contribution to the cross-spectrum covariance is:
\begin{equation}
\mathrm{Cov}^{\mathrm{SSC}}(\Cl^{AB}_{\ell_1},\, \Cl^{CD}_{\ell_2})
= \int \frac{\mathrm{d}\chi}{\chi^4}\;
\frac{\partial C_{\ell_1}^{AB}}{\partial \delta_b(\chi)}\;
\frac{\partial C_{\ell_2}^{CD}}{\partial \delta_b(\chi)}\;
\sigma_b^2(\chi),
\label{eq:ssc}
\end{equation}
where $\sigma_b^2$ is the variance of the background density mode, determined by the survey window function. Crucially, \textbf{SSC does not scale simply as $1/\fsky$}---smaller surveys probe more power in modes just outside their boundary.

For a \textbf{MeerKAT-scale survey} ($\sim 4{,}000$~deg$^2$, $\fsky \approx 0.10$), SSC is significantly more important than for \textbf{SKA} ($\sim 25{,}000$--$30{,}000$~deg$^2$, $\fsky \approx 0.6$--$0.7$).

The response of the angular power spectrum to a background overdensity $\delta_b$ can be computed in the halo model:
\begin{equation}
\frac{\partial \Cl^{ij}}{\partial \delta_b}
= \int \frac{\mathrm{d}\chi}{\chi^2}\; W_i(\chi)\, W_j(\chi)\;
\frac{\partial P_{ij}(k{=}\ell/\chi,\, z)}{\partial \delta_b},
\end{equation}
where the power spectrum response includes the effects of the background mode on halo abundances (via the peak-background split), halo bias, and halo profiles.


\subsection{The 1-Halo Trispectrum and Dark Matter Annihilation}

The 1-halo trispectrum is especially important for DM annihilation signals because the gamma-ray emissivity scales as $\rho_{\mathrm{DM}}^2$, making the power spectrum effectively a two-point function of $\rho^2$ and the covariance an eight-point function of~$\rho$. In the halo model, the 1-halo trispectrum of the $\rho^2$ field is schematically:
\begin{equation}
T^{\mathrm{1h}}_{\delta^2}(k_1, k_2, k_3, k_4) \propto \int \mathrm{d}M\; \dndm\;
\frac{\vtilde(k_1|M)\,\vtilde(k_2|M)\,\vtilde(k_3|M)\,\vtilde(k_4|M)}
{\langle\rho^2\rangle^4},
\end{equation}
where $\vtilde(k|M) = \int \mathrm{d}^3x\;\rho^2(\mathbf{x}|M)\,e^{i\mathbf{k}\cdot\mathbf{x}}$ is the Fourier transform of the squared NFW profile (following the notation of Part~3 of the companion formalism document). Since $\vtilde(0|M) = \int \mathrm{d}^3x\;\rho^2 \propto c^3 M^2 / r_s^3$, the steep dependence on halo concentration makes this term heavily weight concentrated, low-mass halos.


\subsection{Quantitative Impact on HI $\times$ Gamma-Ray Forecasts}

For the specific multipole range $\ell \sim 10$--$2000$ relevant to HI $\times$ gamma-ray cross-correlations, several factors suppress non-Gaussian corrections:
\begin{itemize}[leftmargin=2em]
\item \textbf{Thermal noise} in HI maps dominates the auto-spectrum for MeerKAT (and partly for SKA1), acting as an effective regularizer.
\item \textbf{Photon noise} in Fermi-LAT maps similarly dominates at all but the lowest energies.
\item The covariance is proportional to the product of noise-dominated auto-spectra, dwarfing the non-Gaussian contribution.
\item Existing Fermi-LAT $\times$ galaxy/lensing cross-correlations at similar multipole ranges are well-described by Gaussian covariance (Tr\"oster et al.\ 2017; Ammazzalorso et al.\ 2020).
\end{itemize}

Quantitatively, non-Gaussian corrections are estimated to degrade the SNR by $\lesssim 10\%$ for current noise levels. The exception is SSC for the smaller MeerKAT survey at $\ell \lesssim$ few hundred, which could affect constraints at the $5$--$15\%$ level in the signal-dominated regime. For weak lensing surveys (a useful benchmark), Takada \& Jain (2009) found non-Gaussian errors degrade the cumulative power spectrum S/N by up to a factor of 2 at $\ell \sim 1000$--$3000$, but marginalized parameter constraints degrade by less than 10--20\%. Adding tomographic redshift bins further reduces the non-Gaussian impact by averaging over independent volumes.


%======================================================================
\section{Practical Power Spectrum Estimation}
\label{sec:practical}
%======================================================================

\subsection{The Pseudo-$\Cl$ Estimator and Mode-Coupling Matrix}

On a masked sky, the spherical harmonic transform of the masked field $\tilde{f}(\nhat) = W(\nhat)\, f(\nhat)$ produces pseudo-multipole coefficients $\tilde{a}_{\ell m}$. The pseudo-$\Cl$ estimator
\begin{equation}
\tilde{C}_\ell = \frac{1}{2\ell+1}\sum_m \tilde{a}_{\ell m}^i\, \tilde{a}_{\ell m}^{j*}
\end{equation}
is biased by mode-coupling:
\begin{equation}
\langle \tilde{C}_\ell \rangle = \sum_{\ell'} M_{\ell\ell'}\, C_{\ell'},
\label{eq:mode_coupling_bias}
\end{equation}
where the mode-coupling matrix depends only on the \textbf{mask power spectrum}:
\begin{equation}
M_{\ell\ell'} = \frac{2\ell'+1}{4\pi}
\sum_{\ell''} (2\ell''+1)\, W_{\ell''}
\begin{pmatrix} \ell & \ell' & \ell'' \\ 0 & 0 & 0 \end{pmatrix}^2.
\label{eq:mode_coupling_matrix}
\end{equation}
Here $W_\ell$ is the angular power spectrum of the mask, and the parenthesized term is the Wigner $3j$-symbol enforcing the triangle inequality $|\ell - \ell'| \leq \ell'' \leq \ell + \ell'$.

The \textbf{MASTER algorithm} (Hivon et al.\ 2002) recovers unbiased estimates via:
\begin{equation}
\hat{C}_\ell = \sum_{\ell'} M_{\ell\ell'}^{-1}
\left[\tilde{C}_{\ell'} - \langle N_{\ell'}\rangle\right],
\label{eq:master}
\end{equation}
with additional corrections for beam ($B_\ell^2$), pixel window ($p_\ell^2$), and transfer functions.


\subsection{NaMaster for Cross-Correlations with Different Masks}

\texttt{NaMaster} (Alonso, Sanchez \& Slosar 2019) provides a public, validated implementation of the pseudo-$\Cl$ framework for arbitrary spin-0 and spin-2 fields. For HI $\times$ gamma-ray cross-correlations with different masks $W_{\mathrm{HI}}$ and $W_\gamma$, the \textbf{effective mask is the product} $W_{\mathrm{eff}} = W_{\mathrm{HI}} \times W_\gamma$, and the mode-coupling matrix is computed from its power spectrum. The workflow is:
\begin{enumerate}[leftmargin=2em]
\item Create \texttt{NmtField} objects for each map with their respective masks.
\item Compute the coupling matrix workspace (\texttt{NmtWorkspace}), which encodes $M_{\ell\ell'}$ and can be saved for reuse.
\item Compute the pseudo-$\Cl$ from the masked fields.
\item Decouple into bandpowers using the precomputed workspace.
\end{enumerate}
\texttt{NaMaster} also implements the Gaussian covariance estimation (Garc\'ia-Garc\'ia et al.\ 2019) with $\mathcal{O}(\ell_{\max}^3)$ scaling, and provides \textbf{bandpower window functions} $\mathcal{W}_{b\ell}$ that encode residual mode-coupling after binning---theory predictions must be convolved with these before comparison to data.


\subsection{HEALPix Pixelization Choices}

The appropriate $N_{\mathrm{side}}$ must satisfy $\theta_{\mathrm{pix}} \approx 3438'/N_{\mathrm{side}}$ being 2--3$\times$ smaller than the instrument angular resolution:

\begin{table}[h]
\centering
\begin{tabular}{@{}llcc@{}}
\toprule
Instrument/Energy & Angular resolution & $N_{\mathrm{side}}$ & $\ell_{\max}$ \\
\midrule
Fermi-LAT, $E < 1$~GeV & $\sim 1^\circ$--$3^\circ$ & 64--128 & $\sim 190$--380 \\
Fermi-LAT, 1--10~GeV & $\sim 0.2^\circ$--$0.8^\circ$ & 128--256 & $\sim 380$--770 \\
Fermi-LAT, $E > 10$~GeV & $< 0.2^\circ$ & 512--1024 & $\sim 1500$--3000 \\
MeerKAT single-dish & $\sim 1^\circ$ at L-band & 256--512 & $\sim 770$--1500 \\
SKA interferometer & arcminute-scale & 1024--2048 & $\sim 3000$--6000 \\
\bottomrule
\end{tabular}
\caption{Recommended HEALPix parameters for each instrument and energy range.}
\label{tab:healpix}
\end{table}

The pixel window function $p_\ell^2$ (available via \texttt{healpy.pixwin}) must be corrected for in all analyses. Retaining multipoles up to $\ell_{\max} = 3\,N_{\mathrm{side}} - 1$ avoids aliasing.


\subsection{The Noise-Free Cross-Correlation Advantage}

\textbf{Cross-correlation between HI and gamma-ray maps is automatically free from noise bias} because the noise sources are physically independent. The cross-spectrum estimator satisfies
\begin{equation}
\langle \hat{C}_\ell^{\mathrm{cross}}\rangle = \Cl^{\mathrm{signal}}
\end{equation}
without any noise subtraction---the most computationally expensive and error-prone step in auto-spectrum analyses is entirely eliminated. No noise model, no Monte Carlo noise estimation, and no noise power spectrum subtraction is needed.

For the auto-spectrum, in contrast, one must compute and subtract $\langle N_\ell \rangle$ with sufficient accuracy that the residual $|\hat{C}_\ell^{\mathrm{auto}} - N_\ell - \Cl^{\mathrm{signal}}| \ll \Cl^{\mathrm{signal}}$. This is notoriously difficult for intensity mapping, where thermal noise exceeds the signal by orders of magnitude, and for Fermi-LAT, where the photon noise must be precisely characterized.


\subsection{Empirical Error Estimation}

\paragraph{Jackknife resampling.}
Divide the survey into $N_{\mathrm{JK}}$ spatial patches and compute the power spectrum with each patch removed. The jackknife covariance is:
\begin{equation}
\mathrm{Cov}_{\mathrm{JK}} = \frac{N_{\mathrm{JK}}-1}{N_{\mathrm{JK}}}
\sum_{i=1}^{N_{\mathrm{JK}}}
\left(\hat{\mathbf{C}}^{(i)} - \bar{\mathbf{C}}\right)
\left(\hat{\mathbf{C}}^{(i)} - \bar{\mathbf{C}}\right)^\top.
\label{eq:jackknife}
\end{equation}
This is model-independent and captures non-Gaussian contributions, but tends to \textbf{overestimate errors by 25--60\%}. When inverting the covariance, the Hartlap correction factor $(N_{\mathrm{JK}} - N_{\mathrm{bins}} - 2)/(N_{\mathrm{JK}} - 1)$ is required to obtain an unbiased inverse.

\paragraph{Monte Carlo simulations.}
Generate Gaussian random fields with the measured (or theoretical) power spectrum, convolve with beams, add realistic noise, apply masks, and run the full analysis pipeline. The sample covariance from $N_{\mathrm{sim}} \gg N_{\mathrm{bins}}$ realizations (typically 500--1000) captures all mask effects precisely. For cross-correlations, correlated pairs of maps with the correct cross-spectrum must be generated using the Cholesky decomposition of the $2 \times 2$ spectral matrix.

\paragraph{Best practice.}
Combine \texttt{NaMaster}'s analytical Gaussian covariance as the primary estimate, validated against a limited set of Monte Carlo simulations ($\sim$100--500), with jackknife as an independent model-free cross-check.


\subsection{Null Tests}

Essential null tests to validate the cross-correlation measurement include:
\begin{enumerate}[leftmargin=2em]
\item \textbf{Map rotation}: Rotating one map in Galactic longitude (e.g., by $90^\circ$, $180^\circ$, $270^\circ$) before cross-correlating, destroying true correlations while preserving individual map properties.
\item \textbf{Random catalog cross-correlation}: Replacing one field with randoms matching the survey geometry.
\item \textbf{Temporal splits}: Cross-correlating first-half vs.\ second-half observations to detect time-dependent systematics.
\item \textbf{Energy/frequency consistency}: Verifying the expected spectral dependence of the cross-correlation signal across gamma-ray energy and HI frequency bands.
\item \textbf{$\chi^2$ against null}: Computing $\chi^2_{\mathrm{null}} = \hat{\mathbf{C}}^\top \mathrm{Cov}^{-1} \hat{\mathbf{C}}$ for the measured cross-spectrum and comparing to the $\chi^2(N_{\mathrm{bins}})$ distribution.
\end{enumerate}


%======================================================================
\section{Statistical Outputs and Deliverables}
\label{sec:outputs}
%======================================================================

\subsection{Predicted SNR for Astrophysical Detection}

Pinetti et al.\ Table~4 forecasts the cumulative signal-to-noise ratio for the astrophysical cross-correlation signal:

\begin{table}[h]
\centering
\begin{tabular}{@{}lc@{}}
\toprule
Configuration & SNR \\
\midrule
Fermi-LAT $\times$ MeerKAT & 3.7 \\
Fermi-LAT $\times$ SKA Phase~1 & 5.7 \\
Fermi-LAT $\times$ SKA Phase~2 & 8.2 \\
\bottomrule
\end{tabular}
\caption{Forecast signal-to-noise ratios for the astrophysical HI $\times$ gamma-ray cross-correlation (reproduced from Pinetti et al.\ 2020, Table~4).}
\label{tab:snr}
\end{table}

MeerKAT provides a $\sim 3.7\sigma$ indication, while SKA1 crosses the $5\sigma$ detection threshold and SKA2 achieves a robust detection.


\subsection{Dark Matter Exclusion Limits}

The $2\sigma$ (95.45\% CL) upper limits on $\sigmav$ versus $\mchi$ (via the $\Delta\chi^2$ raster scan with $\Delta\chi^2 = 4$ at 1~dof) show that Fermi-LAT $\times$ SKA can probe the thermal relic cross-section up to $\mchi \approx 130$--$150$~GeV for the $b\bar{b}$ annihilation channel. The DM signal template uses the halo model framework with NFW profiles, incorporating both one-halo and two-halo terms, with the characteristic $\rho_{\mathrm{DM}}^2$ redshift window function (Eq.~4.1 of Pinetti et al.):
\begin{equation}
W_\gamma^{\mathrm{DM}}(E,\chi) = \frac{(\Omega_{\mathrm{DM}}\rho_c)^2}{4\pi}\;
\frac{\sigmav}{2\mchi^2}\;
\frac{(1+z)^3}{H(z)}\;\Delta^2(z)\;
\int_{E_{\min}}^{E_{\max}} \mathrm{d}E_0\;
\left.\frac{\mathrm{d}N_\gamma}{\mathrm{d}E'}\right|_{E'=E_0(1+z)}\;
e^{-\tau(E_0(1+z),z)}.
\end{equation}


\subsection{Real Data Analysis Products}

For actual measurements, the complete analysis chain produces:
\begin{enumerate}[leftmargin=2em]
\item Measured bandpowers $\hat{C}_b$ with error bars from the diagonal of the covariance matrix.
\item Null tests establishing the reality of the cross-correlation signal.
\item Best-fit amplitudes for astrophysical source templates via Eq.~(\ref{eq:wls}).
\item Upper limits or detection of a DM component through template fitting.
\end{enumerate}

The \textbf{spectral template analysis} exploits the energy dependence of the cross-power spectrum to distinguish source populations---blazars dominate at higher energies with hard power-law spectra, while DM templates exhibit a characteristic spectral cutoff at $E = \mchi$ that shifts with assumed WIMP mass.


\subsection{Frequentist Versus Bayesian Reporting}

\textbf{Confidence intervals} (frequentist) state that 95\% of identically constructed intervals contain the true parameter---a statement about procedure reliability. \textbf{Credible intervals} (Bayesian) state that the parameter has 95\% posterior probability of lying within the interval---a direct probability statement that depends on the prior.

For DM cross-section limits, the \textbf{Feldman--Cousins unified approach} prevents pathologies (empty intervals, flip-flopping between limits and detections) by using a likelihood-ratio ordering principle that automatically transitions between upper limits and two-sided intervals based on the data.


\subsection{The Look-Elsewhere Effect}

The look-elsewhere effect from scanning over $\mchi$ (and potentially annihilation channels) inflates local $p$-values. Following Gross \& Vitells (2010), the global significance is approximately:
\begin{equation}
p_{\mathrm{global}} \approx p_{\mathrm{local}} + \langle N_{\mathrm{up}} \rangle \cdot P(\chi^2_1 > u),
\end{equation}
where $\langle N_{\mathrm{up}} \rangle$ is the expected number of upcrossings of the test statistic above the local threshold~$u$, related to the Euler characteristic of the parameter space.

The Bayesian framework handles this naturally through prior-weighted marginalization over $\mchi$, effectively paying the ``trials factor'' through the prior volume. A Savage--Dickey density ratio can directly assess whether the data prefer $\sigmav = 0$ (null) versus $\sigmav > 0$ (alternative).


%======================================================================
\section{Key References}
\label{sec:references}
%======================================================================

\subsection{Primary Papers}

\begin{itemize}[leftmargin=2em]
\item Pinetti, Camera, Fornengo \& Regis (2020), JCAP 07, 044. arXiv:2001.02649. \textit{The foundational paper for HI $\times$ gamma-ray cross-correlation DM searches.}
\item Fornengo \& Regis (2014), Front.\ Physics 2, 6. \textit{Halo model formalism for density $\times$ density-squared cross-correlations.}
\item Camera, Fornasa, Fornengo \& Regis (2015), JCAP 06, 029. arXiv:1411.4651. \textit{Tomographic-spectral approach for DM in gamma-ray $\times$ LSS.}
\end{itemize}

\subsection{Statistical Methodology}

\begin{itemize}[leftmargin=2em]
\item Knox (1995), Phys.\ Rev.\ D 52, 4307. \textit{Original derivation of the power spectrum variance formula.}
\item Hivon et al.\ (2002), ApJ 567, 2. \textit{MASTER pseudo-$\Cl$ algorithm.}
\item Alonso, Sanchez \& Slosar (2019), MNRAS 484, 4127. \textit{NaMaster: unified pseudo-$\Cl$ framework.}
\item Cowan, Cranmer, Gross \& Vitells (2011), Eur.\ Phys.\ J.\ C 71, 1554. \textit{Asymptotic formulae for profile likelihood tests.}
\item Feldman \& Cousins (1998), Phys.\ Rev.\ D 57, 3873. \textit{Unified confidence intervals.}
\item Gross \& Vitells (2010), Eur.\ Phys.\ J.\ C 70, 525. \textit{Look-elsewhere effect.}
\item Trotta (2008), Contemp.\ Phys.\ 49, 71. \textit{Bayesian model comparison in cosmology.}
\end{itemize}

\subsection{Non-Gaussian Covariance}

\begin{itemize}[leftmargin=2em]
\item Takada \& Hu (2013), Phys.\ Rev.\ D 87, 123504. arXiv:1302.6994. \textit{Power spectrum super-sample covariance.}
\item Lacasa \& Grain (2019), A\&A 615, A1. \textit{Halo model covariance of the galaxy angular power spectrum.}
\item Barreira, Krause \& Schmidt (2018), JCAP 10, 053. \textit{Complete treatment of SSC for galaxy clustering and lensing.}
\end{itemize}

\subsection{Gamma-Ray Cross-Correlation Analyses}

\begin{itemize}[leftmargin=2em]
\item Cuoco, Xia, Regis, Branchini, Fornengo \& Viel (2015), ApJS 221, 29. arXiv:1506.01030. \textit{DM searches via UGRB $\times$ galaxy catalogs.}
\item Ammazzalorso et al.\ (2020), Phys.\ Rev.\ Lett.\ 124, 101102. \textit{DES weak lensing $\times$ Fermi-LAT cross-correlation.}
\item Zhou, Bernal, Pinetti et al.\ (2024), arXiv:2410.00375. \textit{UGRB $\times$ DESI cross-correlation forecasts.}
\item Ackermann et al.\ (2018), Phys.\ Rev.\ D 98, 022002. \textit{Fermi-LAT UGRB auto-correlation measurement.}
\end{itemize}

\subsection{Software Packages}

\begin{itemize}[leftmargin=2em]
\item \texttt{NaMaster}: \url{https://github.com/LSSTDESC/NaMaster}
\item \texttt{healpy}/HEALPix: \url{https://healpy.readthedocs.io}
\item \texttt{Fermitools}: \url{https://fermi.gsfc.nasa.gov/ssc/data/analysis/software/}
\item PPPC4DMID: \url{http://www.marcocirelli.net/PPPC4DMID.html}
\item \texttt{ebltable}: \url{https://github.com/me-manu/ebltable}
\item \texttt{emcee}: \url{https://emcee.readthedocs.io}
\item \texttt{MultiNest}/\texttt{PyMultiNest}: \url{https://github.com/JohannesBuchner/PyMultiNest}
\item \texttt{Cobaya}: \url{https://cobaya.readthedocs.io}
\item \texttt{PolyChord}: \url{https://github.com/PolyChord/PolyChordLite}
\end{itemize}


%======================================================================
\section{Concluding Remarks}
\label{sec:conclusion}
%======================================================================

The statistical framework for HI $\times$ gamma-ray cross-correlation analysis rests on a remarkably clean theoretical foundation. The Gaussian covariance derived from Wick's theorem through the Knox formula provides an accurate, analytically tractable error model for most of the parameter space relevant to MeerKAT and SKA observations. The \textbf{noise-free nature of cross-correlations}---perhaps the single most important statistical advantage---eliminates the need for noise modeling entirely, while the tomographic capability of intensity mapping provides the spectral-spatial discrimination needed to separate dark matter from astrophysical backgrounds.

Three insights emerge that are not immediately obvious from the standard treatment:

\begin{enumerate}[leftmargin=2em]
\item The interplay between gamma-ray PSF degradation at low energies and photon statistics improving at low energies creates an \textbf{energy-dependent sweet spot} that couples to the multipole-space sweet spot. Optimal sensitivity comes from carefully weighting the two-dimensional $(E, \ell)$ information space.

\item \textbf{Non-Gaussian covariance is likely subdominant} ($\lesssim 10\%$ effect on SNR) for current instruments precisely because noise dominance in both auto-spectra regularizes the covariance, though SSC warrants monitoring for the smaller MeerKAT survey area.

\item The transition from simple $\Delta\chi^2$ to profile likelihood or Bayesian inference is not merely a theoretical refinement---it becomes essential when marginalizing over uncertain astrophysical models that share spectral and angular features with the DM signal, particularly for source classes like star-forming galaxies whose redshift evolution partially mimics the DM window function.
\end{enumerate}

\end{document}
